\begin{figure}[p]
\centering
\begin{tikzpicture}[
  font=\footnotesize,
  node distance=6mm and 8mm,
  box/.style={
    draw,
    rounded corners,
    align=center,
    inner sep=4pt,
    minimum width=30mm,
    minimum height=8mm
  },
  arr/.style={-Latex, line width=0.6pt}
]

% ===== 第一行：数据与规则构建 =====
\node[box] (data)
{真实多模态样本\\
文本问题+图像};

\node[box, right=of data] (annot)
{人工复核与\\
类别标注};

\node[box, right=of annot] (principle)
{类别映射为\\
安全原则};

\node[box, right=of principle] (rule)
{结合具体问题\\
生成实例规则};

% ===== 第二行：两阶段合成 =====
\node[box, below=of principle] (cot)
{阶段1：生成安全\\
推理过程};

\node[box, right=of cot] (answer)
{阶段2：生成\\
最终回复};

\node[box, below=of answer] (output)
{结构化训练样本\\
规则+推理+回复};

% ===== 连接线 =====
\draw[arr] (data) -- (annot);
\draw[arr] (annot) -- (principle);
\draw[arr] (principle) -- (rule);

\draw[arr] (rule) -- (cot);
\draw[arr] (cot) -- (answer);
\draw[arr] (answer) -- (output);

\end{tikzpicture}
\caption{基于规则的两阶段安全数据合成流程。类别首先映射为抽象安全原则，\\
再结合具体问题生成实例级规则；模型先进行风险分析，再在规则约束下生成最终回复。}
\label{fig:data_synthesis_flow}
\end{figure}